\documentclass[12pt]{article}
\usepackage[T1]{fontenc}
\usepackage[utf8]{inputenc}
\usepackage{lmodern}
\usepackage{textcomp}
\usepackage{enumitem}
\usepackage{geometry}
\usepackage{graphicx}
\usepackage{amsmath, amssymb}
\geometry{margin=1in}
\begin{document}
\begin{center}
Constitutional Law - Undergraduate Level Assessment 10 \
Total Marks: 40 \
Answer all questions.
\end{center}

\section*{Multiple Choice (10 marks)}
\begin{enumerate}[label=\arabic*.]
\item Which one of the following statements is correct?
\begin{enumerate}[label=(\alph*)]
    \item The Council of Europe was established in 1950 and consists of 27 member states
    \item The Council of Europe was established in 1949 and consists of 47 member states
    \item The Council of Europe was established in 1959 and consists of 34 member states
    \item The Council of Europe was established in 1984 and consists of 19 member states
\end{enumerate}
\item Iverson Jewelers wrote a letter to Miller, 'We have received an exceptionally fine self winding Rolox watch which we will sell to you at a very favorable price.'
\begin{enumerate}[label=(\alph*)]
    \item The letter is an offer to sell
    \item A valid offer cannot be made by letter.
    \item The letter contains a valid offer which will terminate within a reasonable time.
    \item The letter lacks one of the essential elements of an offer.
\end{enumerate}
\item Was the use of armed force permitted prior to the United Nations Charter?
\begin{enumerate}[label=(\alph*)]
    \item Armed force was prohibited
    \item Armed force was permitted with no restrictions
    \item Armed force was permitted subject to few restrictions
    \item Armed force was not regulated under international law prior to 1945
\end{enumerate}
\item What types of force does Article 2(4) of the UN Charter prohibit?
\begin{enumerate}[label=(\alph*)]
    \item Article 2(4) encompasses only armed force
    \item Article 2(4) encompasses all types of force, including sanctions
    \item Article 2(4) encompasses all interference in the domestic affairs of States
    \item Article 2(4) encompasses force directed only against a State's territorial integrity
\end{enumerate}
\item Which of the following is the strongest argument against ethical relativism's hostility to human rights?
\begin{enumerate}[label=(\alph*)]
    \item Utilitarianism
    \item Communitarianism.
    \item Cognitivism.
    \item Positivism.
\end{enumerate}
\end{enumerate}

\section*{True / False (10 marks)}
\textit{Answer True or False}
\begin{enumerate}[label=\arabic*.]
\setcounter{enumi}{5}
\item True or False: The principle that punishment should fit the crime is inconsistent with consequentialist justifications of punishment.
\item True or False: All types of human rights are indivisible, interrelated, and interdependent, as stated in the UN Vienna Declaration of 1993.
\item True or False: Normative legal theory is solely concerned with the strict application of laws without regard to moral implications.
\item True or False: Radical feminists criticize liberal feminists for promoting the idea that achieving equality means making women similar to men.
\item True or False: Many Critical Race Theorists view concepts like 'justice' as irrelevant to contemporary social issues faced by people of colour.
\end{enumerate}

\section*{Long Form Response (20 marks)}
\textit{Answer each question with a clear and complete explanation.}
\begin{enumerate}[label=\arabic*.]
\setcounter{enumi}{10}
\item Discuss the implications of the Kadi judgment on the incorporation of UN Security Council resolutions, particularly in relation to human rights standards.
\item Discuss the circumstances under which armed violence by non-State actors can be considered an armed attack under Article 51 of the UN Charter, referencing the Caroline case as a precedent.
\end{enumerate}

\vfill
\noindent\textit{End of Paper}
\end{document}